\documentclass[a4paper]{article}

% --- Основные пакеты ---
\usepackage{graphicx}     % картинки
\usepackage{geometry}     % поля страницы
\usepackage{amsmath}      % математика
\usepackage{iftex}        % определение движка (pdfLaTeX / XeLaTeX / LuaLaTeX)
\ifPDFTeX
    \usepackage{amssymb}       % математические символы
    \usepackage[T2A]{fontenc}  % кириллица
    \usepackage[utf8]{inputenc}% кодировка UTF-8
    \usepackage{anyfontsize}   % точные размеры шрифтов без подстановки
\else
    \usepackage{unicode-math}  % масштабируемые OpenType-матшрифты
    \usepackage{fontspec}      % системные шрифты с поддержкой кириллицы
    \setmainfont{Times New Roman}
    \setsansfont{Arial}
    \setmonofont{Consolas}
\fi
\usepackage[russian]{babel} % многоязычность
\usepackage{caption}      % настройка подписей
\usepackage[fontsize=13pt]{fontsize} % Размер шрифта
\usepackage{float}        % [H] позиционирование фигур
\usepackage{xcolor}       % цвета (для листингов)
\usepackage{listings}     % листинги исходного кода
\usepackage{hyperref}     % гиперссылки (подключаем одним из последних)

% --- Настройка страницы ---
\geometry{top=1.5cm, bottom=1.5cm, left=3cm, right=1.5cm}

% --- Настройка гиперссылок ---
\hypersetup{
    colorlinks,
    linkcolor=black,
    urlcolor=blue,
    citecolor=black
}

% --- Настройка подписей ---
\renewcommand{\figurename}{Рис.}
\captionsetup{labelsep=period}

% --- Настройка листингов ---
\definecolor{lstback}{RGB}{248,248,248}
\definecolor{lstcomment}{RGB}{0,128,0}
\definecolor{lstkeyword}{RGB}{0,0,180}
\definecolor{lststring}{RGB}{163,21,21}
\definecolor{lstnumber}{RGB}{128,128,128}

\lstset{
    language=C++,
    basicstyle=\ttfamily\footnotesize,
    keywordstyle=\color{lstkeyword}\bfseries,
    commentstyle=\color{lstcomment}\itshape,
    stringstyle=\color{lststring},
    numberstyle=\tiny\color{lstnumber},
    numbers=left,
    numbersep=8pt,
    frame=single,
    rulecolor=\color{gray!50},
    backgroundcolor=\color{lstback},
    breaklines=true,
    breakatwhitespace=false,
    showstringspaces=false,
    tabsize=4,
    columns=fullflexible,
    keepspaces=true,
    extendedchars=true,
    inputencoding=utf8,
    captionpos=b
}
\renewcommand{\lstlistingname}{Листинг}

% Заголовки разделов и подразделов без нумерации
\newcommand{\mainsection}[1]{%
    \phantomsection
    \section*{\centering\fontsize{16pt}{0pt}\selectfont\textbf{#1}}%
}
\newcommand{\subsectiontext}[1]{%
    \section*{\fontsize{13pt}{0pt}\selectfont\textbf{#1}}%
}

\begin{document}

% --- Титульная страница ---
\centering
Министерство науки и высшего образования Российской Федерации\\
Федеральное государственное автономное образовательное учреждение\\
высшего образования\\
«Южно-Уральский государственный университет\\
(национальный исследовательский университет)»\\
Институт естественных и точных наук

\setlength{\parskip}{7.5cm}

Отчет\\
Лабораторная работа №2.\\
«Броуновское движение частиц»

\setlength{\parskip}{5.9cm}

\raggedleft
\begin{minipage}{0.46\textwidth}
\raggedright
Выполнил: студент группы ЕТ-310\\
Тельнов С.М.\\
Проверил: ассистент кафедры ВМ\\
Помыкалов Е.В.
\end{minipage}

\setlength{\parskip}{5.9cm}

\centering
Челябинск — 2026

\setlength{\parskip}{0pt plus 1pt}
\newpage

% --- Содержание ---
\centering
\textbf{Краткое содержание}
\begin{itemize}
    \setlength{\parskip}{0.4cm}
    \item \hyperref[sec:theory]{\text{Теоретическая часть} \dotfill \pageref{sec:theory}}
    \item \hyperref[sec:practice]{\text{Практическая часть} \dotfill \pageref{sec:practice}}
    \item \hyperref[sec:conclusion]{\text{Вывод} \dotfill \pageref{sec:conclusion}}
\end{itemize}
\newpage

\setlength{\parindent}{1.5em}

\raggedright

% --- Теоретическая часть ---
\mainsection{Теоретическая часть}
\label{sec:theory}

В лабораторной работе моделируется плоское движение ансамбля частиц (шаров) в замкнутой прямоугольной области. Частицы движутся равномерно и прямолинейно (сила тяжести отсутствует) и испытывают абсолютно упругие соударения друг с другом и со стенками, благодаря чему суммарная кинетическая энергия системы сохраняется. Визуализация процесса выполняется в реальном времени средствами библиотеки SFML, а управление сведено к минимуму: создание частицы перетаскиванием мыши (вектор натяжения задаёт начальную скорость), изменение радиуса частиц колесом мыши, пауза пробелом и кнопка генерации случайных частиц.

\subsectiontext{Библиотека SFML}

\textbf{SFML} (Simple and Fast Multimedia Library) --- это кроссплатформенная объектно-ориентированная библиотека для языка C++, предоставляющая удобный интерфейс к графике, звуку, работе с окнами и сетью. Библиотека разделена на несколько независимых модулей:

\begin{itemize}
    \item \texttt{System} --- базовые утилиты: работа со временем (\texttt{sf::Clock}, \texttt{sf::sleep}), потоками, математические типы \texttt{sf::Vector2f}, \texttt{sf::Vector2i};
    \item \texttt{Window} --- создание окна \texttt{sf::RenderWindow}, обработка событий \texttt{sf::Event} (клавиатура, мышь, закрытие окна), ввод;
    \item \texttt{Graphics} --- рисование графических примитивов (\texttt{sf::CircleShape}, \texttt{sf::RectangleShape}, \texttt{sf::Text}, \texttt{sf::Vertex}), работа с текстурами, шрифтами \texttt{sf::Font} и источниками отрисовки;
    \item \texttt{Audio} и \texttt{Network} --- воспроизведение звука и сетевые соединения (в данной работе не используются).
\end{itemize}

Библиотека является многопоточной, поддерживает OpenGL и работает на Windows, Linux и macOS. В настоящей работе используются три модуля: \texttt{System}, \texttt{Window} и \texttt{Graphics}.

Проект собирается системой сборки CMake (целевой стандарт --- C++17). Поиск библиотеки осуществляется командой \texttt{find\_package(SFML 3 COMPONENTS Graphics Window System REQUIRED)}, а полученные цели \texttt{SFML::Graphics}, \texttt{SFML::Window}, \texttt{SFML::System} линкуются к исполняемому файлу \texttt{program}. Шрифт \texttt{DejaVuSans.ttf} копируется в каталог сборки, так как он необходим для вывода текста на экран.

Основная идея архитектуры приложения --- классический игровой цикл: на каждой итерации сначала опрашиваются и обрабатываются события окна (\texttt{window.pollEvent}), затем выполняется шаг физического моделирования, после чего сцена полностью перерисовывается. Такое разделение удобно, поскольку обработка ввода, расчёт физики и вывод изображения независимы друг от друга.

\subsectiontext{Модель движения частиц}

Каждая частица описывается структурой \texttt{Particle}, хранящей координаты центра $(x, y)$, проекции скорости $(v_x, v_y)$ и массу $m$. Ускорения отсутствуют, поэтому движение между соударениями --- равномерное и прямолинейное. Интегрирование уравнений движения выполняется явным методом Эйлера первого порядка с постоянным шагом $\Delta t$:

\begin{equation}
    x_i^{\,n+1} = x_i^{\,n} + v_{x,i}^{\,n}\,\Delta t, \qquad
    y_i^{\,n+1} = y_i^{\,n} + v_{y,i}^{\,n}\,\Delta t.
\end{equation}

Шаг интегрирования в программе равен $\Delta t = 0{,}01$~с. Суммарная кинетическая энергия системы вычисляется как

\begin{equation}
    E = \sum_{i=1}^{N} \frac{m_i}{2}\left(v_{x,i}^2 + v_{y,i}^2\right),
\end{equation}

где $N$ --- число частиц. Именно эта величина контролируется на экране: для физически корректной модели она должна оставаться постоянной, а отклонение текущей энергии от начальной показывает накопленную численную ошибку.

Отскок от стенок моделируется упруго: если центр частицы вышел за допустимую границу (учитывается радиус $r$ и толщина рамки), координата возвращается на границу, а соответствующая проекция скорости меняет знак, например:

\begin{equation}
    x_i < r \;\Rightarrow\; x_i = r,\qquad v_{x,i} = -v_{x,i}.
\end{equation}

\subsectiontext{Упругое соударение двух шаров}

Наиболее содержательная часть модели --- обработка парных соударений шаров. Рассмотрим две частицы 1 и 2 с массами $m_1$, $m_2$, положениями $\vec{r}_1$, $\vec{r}_2$ и скоростями $\vec{v}_1$, $\vec{v}_2$. Введём вектор взаимного положения и относительную скорость:

\begin{equation}
    \vec{d} = \vec{r}_1 - \vec{r}_2, \qquad d = |\vec{d}|, \qquad \vec{u} = \vec{v}_1 - \vec{v}_2.
\end{equation}

Соударение возможно, только если расстояние между центрами не превышает суммы радиусов, $d \le 2r$ (в программе сравниваются квадраты расстояний, чтобы избежать извлечения корня). Кроме того, пару обрабатывают лишь тогда, когда частицы сближаются, а не расходятся. Это условие записывается через скалярное произведение

\begin{equation}
    \mathrm{dot} = \vec{u}\cdot\vec{d} = (v_{x,1}-v_{x,2})\,d_x + (v_{y,1}-v_{y,2})\,d_y < 0.
\end{equation}

Если $\mathrm{dot} \ge 0$, частицы уже удаляются друг от друга, и обмена импульсами не происходит.

При абсолютно упругом ударе силы взаимодействия направлены вдоль линии центров, поэтому меняются только нормальные к поверхности касания составляющие скоростей, а тангенциальные сохраняются. Из законов сохранения импульса и кинетической энергии для двух тел получается классическая формула для импульса (импульса удара), передаваемого вдоль нормали $\hat{n} = \vec{d}/d$:

\begin{equation}
    J = \frac{2\,m_1 m_2}{m_1 + m_2}\,\bigl(\vec{u}\cdot\hat{n}\bigr).
\end{equation}

Новые скорости находятся как

\begin{equation}
    \vec{v}_1' = \vec{v}_1 - \frac{J}{m_1}\,\hat{n}, \qquad
    \vec{v}_2' = \vec{v}_2 + \frac{J}{m_2}\,\hat{n}.
\end{equation}

Подставляя $\hat{n} = \vec{d}/d$ и $\vec{u}\cdot\hat{n} = \mathrm{dot}/d$, после несложных преобразований получаем формулы, непосредственно реализованные в функции \texttt{stepParticles}:

\begin{equation}
    \vec{v}_1' = \vec{v}_1 - k_1\,\vec{d}, \qquad
    \vec{v}_2' = \vec{v}_2 + k_2\,\vec{d},
    \label{eq:collision}
\end{equation}

где коэффициенты

\begin{equation}
    k_1 = \frac{2\,m_2}{m_1 + m_2}\,\frac{\mathrm{dot}}{d^2}, \qquad
    k_2 = \frac{2\,m_1}{m_1 + m_2}\,\frac{\mathrm{dot}}{d^2}.
\end{equation}

В программе все частицы имеют одинаковую массу $m_1 = m_2 = m$, поэтому формулы упрощаются: $k_1 = k_2 = \mathrm{dot}/d^2$, то есть частицы просто обмениваются нормальными составляющими скоростей.

Из формулы \eqref{eq:collision} видно, что $m_1 \Delta\vec{v}_1 + m_2 \Delta\vec{v}_2 = 0$ --- закон сохранения импульса выполняется точно. Кинетическая энергия также сохраняется, поскольку удар упругий, что и обеспечивает постоянство величины $E$ из выражения (2).

Отдельно рассмотрим проблему взаимного перекрытия. После одной итерации с конечным шагом $\Delta t$ частицы могут «застрять» друг в друге. Чтобы этого избежать, после обновления скоростей производится геометрическая коррекция: частицы раздвигаются вдоль нормали на половину глубины проникновения каждая:

\begin{equation}
    \vec{r}_1 \mathrel{+}= \frac{d_{\min} - d}{2}\,\hat{n}, \qquad
    \vec{r}_2 \mathrel{-}= \frac{d_{\min} - d}{2}\,\hat{n}, \qquad d_{\min} = 2r.
\end{equation}

После сдвига положение каждой частицы ещё раз проверяется на выход за границы области (\texttt{resolveWalls}), чтобы коррекция не вытолкнула шар за стенку.

\subsectiontext{Класс TextBox}

Для удобного создания элементов интерфейса (информационных панелей и кнопок управления) реализован вспомогательный класс \texttt{TextBox}. Он представляет собой прямоугольный блок с рамкой и центрированным текстом и наследуется от \texttt{sf::Drawable}, благодаря чему может отрисовываться тем же вызовом \texttt{window.draw}, что и стандартные примитивы SFML. Исходный код класса приведён в листинге~\ref{lst:textbox}.

Ключевые решения реализации:

\begin{itemize}
    \item шрифт \texttt{sf::Font} хранится \emph{по значению} внутри класса. Объекты \texttt{sf::Text} в SFML хранят лишь указатель на шрифт, поэтому копирование шрифта внутрь виджета гарантирует корректное время жизни и отсутствие висячих ссылок; по этой же причине класс объявлен некопируемым (\texttt{delete} для конструктора копирования и оператора присваивания);
    \item метод \texttt{setPosition} не только задаёт положение прямоугольника, но и пересчитывает начало координат текста (\texttt{setOrigin}) так, чтобы текст всегда был отцентрирован по блоку. Тот же пересчёт выполняется при любом изменении содержимого или размера шрифта, поэтому виджет не «разъезжается» при смене надписи;
    \item метод \texttt{contains} проверяет попадание курсора в границы блока и используется для обработки нажатий на кнопку (а также для подсветки кнопки при наведении);
    \item метод \texttt{draw} просто рисует прямоугольник и текст, что соответствует контракту интерфейса \texttt{sf::Drawable}.
\end{itemize}

Благодаря такой инкапсуляции создание кнопки сводится к нескольким вызовам настройки (\texttt{setSize}, \texttt{setPosition}, \texttt{setString}, цвета), а вся логика центрирования и отрисовки переиспользуется для всех элементов панели.

\begin{lstlisting}[caption={Класс TextBox --- вспомогательный UI-виджет}, label={lst:textbox}]
class TextBox : public sf::Drawable {
    sf::Font font;                      // Шрифт копируется внутрь для корректного времени жизни
    sf::Vector2f size;
    sf::Vector2f position;
    sf::RectangleShape box;
    sf::Text text;

public:
    explicit TextBox(const sf::Font& f) : font(f), text(font, "", 14) {
        box.setOutlineThickness(1.f);
    }

    TextBox(const TextBox&) = delete;
    TextBox& operator=(const TextBox&) = delete;

    void setPosition(sf::Vector2f p) {
        position = p;
        box.setPosition(position);

        const sf::FloatRect bounds = text.getLocalBounds();
        text.setOrigin({bounds.position.x + bounds.size.x / 2.f,
                        bounds.position.y + bounds.size.y / 2.f});
        text.setPosition({position.x + size.x / 2.f, position.y + size.y / 2.f});
    }

    void setSize(sf::Vector2f s) {
        size = s;
        box.setSize(size);
        setPosition(position);
    }

    void setOutlineBoxThickness(float s) { box.setOutlineThickness(s); }

    void setString(const std::string& s) {
        text.setString(s);
        setPosition(position);
    }

    void setCharacterSize(unsigned s) {
        text.setCharacterSize(s);
        setPosition(position);
    }

    void setTextColor(sf::Color c) {
        text.setFillColor(c);
        text.setOutlineColor(c);
    }

    void setFillBoxColor(sf::Color c) { box.setFillColor(c); }
    void setOutlineBoxColor(sf::Color c) { box.setOutlineColor(c); }

    /// Проверка попадания точки в границы блока (для кнопок)
    bool contains(sf::Vector2f p) const {
        return p.x >= position.x && p.x <= position.x + size.x &&
               p.y >= position.y && p.y <= position.y + size.y;
    }

    void draw(sf::RenderTarget& target, sf::RenderStates states) const override {
        target.draw(box, states);
        target.draw(text, states);
    }

    ~TextBox() override = default;
};
\end{lstlisting}

\newpage

% --- Практическая часть ---
\mainsection{Практическая часть}
\label{sec:practice}

Ниже приведён полный исходный код программы (\texttt{src/main.cpp}). Управление: перетаскивание левой кнопкой мыши задаёт начальную скорость новой частицы, колесо мыши изменяет радиус частиц, пробел ставит симуляцию на паузу, кнопка «+10 random balls» добавляет десять случайных частиц. Параметры каждой итерации (номер, число частиц, радиус, кинетическая энергия и её отклонение от начальной) записываются в файл \texttt{simulation\_log.csv}.

\lstinputlisting[caption={Полный исходный код программы (src/main.cpp)}, label={lst:main}]{../src/main.cpp}

\newpage

% --- Вывод ---
\mainsection{Вывод}
\label{sec:conclusion}

В ходе лабораторной работы была разработана программа, моделирующая физически правдоподобное броуновское движение ансамбля частиц в замкнутой области. Движение реализовано методом Эйлера, а взаимодействие шаров друг с другом и со стенками --- абсолютно упругими соударениями, что обеспечивает сохранение суммарной кинетической энергии системы.

Для визуализации процесса использована библиотека SFML, организованная по принципу игрового цикла «события --- физика --- отрисовка». Для удобного построения интерфейса реализован класс \texttt{TextBox}, инкапсулирующий прямоугольный блок с центрированным текстом и поддерживающий обработку нажатий. Управление программой сведено к минимуму кнопок и жестов мыши.

Результаты работы программы контролируются выводом на экран текущей кинетической энергии и её отклонения от начальной: наблюдаемое отклонение остаётся малым, что подтверждает корректность выбранных формул упругого удара и сохранение энергии в модели.

\end{document}
